\section{Literature Positioning}

The concept of enterprise and organizational identity has been widely discussed
in management and organizational studies, predominantly in narrative, cultural,
symbolic, or branding-oriented terms. Within this literature, identity is
typically framed as a shared self-description, a system of values, or a
communicative construct maintained through discourse and sense-making
\cite{albert1985organizational,gioia2000identity,ravasi2016organizational}.

The present work is positioned orthogonally to this tradition. Enterprise
identity is not treated as a descriptive, interpretive, or symbolic phenomenon,
but as a structural property of a formally defined continuum. The defining
elements are admissibility, invariant subspaces, and sustaining cycles within
$\Omega_{\mathrm{EA}}$, rather than narratives, symbols, or stakeholder
perceptions.

Related concepts such as organizational inertia, path dependence, and structural
lock-in address persistence of behavior under constraint
\cite{arthur1989competing,hedberg1981strategy}. However, these approaches do not
distinguish between persistence with and without identity-supporting cycles. As a
result, continued operation is often implicitly equated with identity
persistence. The OC/ETS framework makes this distinction explicit by separating
admissible operation from the existence of invariant identity-supporting
structure.

Literature on organizational failure, decline, and turnaround typically
associates loss of identity with collapse, disintegration, or terminal
dysfunction \cite{weitzel1989decline,hambrick1983some}. By contrast, the present
analysis demonstrates that identity loss can occur as a distinct phase transition
while the enterprise continuum remains admissible and operational. In this
respect, identity loss is shown to be neither synonymous with failure nor
reducible to inefficiency, misalignment, or poor performance.

Finally, while some systems-theoretic and complexity-oriented approaches invoke
notions of structure, stability, and regime shifts
\cite{holland1992complex,scheffer2009critical}, they rarely specify admissibility
boundaries, threshold conditions, or cycle-based invariants in a formal sense.
The contribution of this work is to locate enterprise identity loss precisely
within such a formal phase-space framework, enabling sharp distinctions, no-go
results, and falsifiability conditions that are not available in narrative-based
or metaphorical theories.
